\chapter{Conclusions}
\label{chap:conclusions}

\paragraph{Main findings.} This thesis set out to test a single hypothesis: can natural language serve as a bridge between frozen foundation models? To answer it, we addressed one of the most demanding settings in computer vision, object correspondence across the egocentric and exocentric views of \textit{Ego-Exo4D}. The ``Two Eyes and One Mouth'' pipeline of \cref{chap:pipeline} treats language as that bridge, mapping an object between two viewpoints through a single textual description, with no component trained on the benchmark. This design answers the two limitations of prior work named in \cref{chap:intro} directly. The first is the dependence on training. The second is the difficulty of explaining a correspondence, because earlier methods bury their reasoning in complex fusion modules and abstract learned features. On the official test split, the pipeline reaches within a few IoU points of O-MaMa, the 2025 challenge winner, and overtakes ObjectRelator. It remains the only entry in \cref{tab:main} that is at once free of training, agnostic to the camera, and agnostic to the dataset. On the validation analysis, once its failing cases are set aside, its working mode draws level with even LM-EEC, the strongest published method to date.

The mechanistic analysis of \cref{chap:experiments} does not contradict that hypothesis, but instead sharpens it. The deficit to the state of the art is not a uniform blur of imprecision. It is a discrete ``dead mass'' that holds roughly a third of cases, and these point to optimisations rather than insurmountable failures. The pipeline executes the cross-view correspondence on a few frames and then propagates the mask with a video tracker. When the bridge fails on those frames, nothing downstream can recover or even check the prediction, so the entire instance is lost. The pattern is clear once we break the results down by the quality of the anchor. The pipeline competes with LM-EEC wherever the anchor is good, and cedes almost every case wherever it is poor, which means that the loss is born at the bridge. Two factors push a case into the dead mass, the object's size and the name the VLM writes for it, and the naming is the more damaging of the two. Crucially, the legibility of the language bridge is what turned this gap from a verdict into a diagnosis, marking where nuance is still missing rather than where the paradigm broke down.

Following the diagnosis, we tested a set of targeted interventions. The most obvious repair was to condition the description on a vocabulary of objects plausible for the scene. The gain did not survive a larger vocabulary, however, which suggests that the issue lies in perception rather than in the selection of a name. We then tightened the grounding by passing the detector's box to the segmenter as a crop. This rescued the cases the pipeline was already failing, yet it corrupted the many it already handled, so an unconditional hint hurts on balance and only a gated one could help. We also ruled out a simpler explanation, since raising the brightness of the often dark egocentric frames did not help, which confirms that the dead mass is a failure of naming and not of raw visibility. None of these interventions closed the gap on its own. Each one narrowed where the true fix must act, and that is the dividend of a pipeline open to inspection at every stage (\cref{sec:interventions}).

\paragraph{Future work.} First, the language bridge of \cref{subsec:Stage2} should be constrained to emit a canonical noun phrase while preserving the scene context the name depends on. Naming dominates the dead mass, and the failure is perceptual rather than lexical, since the model nominates the larger neighbour it can actually see once the target shrinks to a fraction of the frame. The obvious remedy of zooming in backfires, because a tight crop destroys the very context that identifies a small object. The constraint must therefore add discipline to the name without removing the scene around it, and the scene vocabulary conditioning discussed above is the most direct attempt at this and the clearest candidate for a full evaluation.

Second, a stronger grounding model would address the slice of the waterfall where the detector boxes a correctly named object off target. LocateAnything \cite{wang2026locateanythingfasthighqualityvisionlanguage}, presented in a 2026 technical report from NVIDIA and not yet peer reviewed, is a natural candidate for this stage.

Third, an external verification gate is needed to catch the confident mis-groundings that make up the rest of the dead mass. These failures are silent, because the pipeline almost always produces a mask, and its own confidence overlaps almost entirely between hits and misses. The lesson of the crop hint applies directly here, namely that an intervention must fire only on the failing population.

Finally, a trigger that reacquires the object after runs of empty masks would recover the propagation drift that surfaces only once the target is widened beyond the exact failures. Unlike the dead mass, this failure announces itself through a long stretch of empty predictions, and is therefore exploitable.

The four directions above refine the pipeline for this task. The paradigm of language as a bridge raises a second tier of questions that outlast it, and three stand out. The first concerns how a VLM is guided to perceive before it names. Because naming dominates the dead mass and the failure is perceptual, the lever is how the model attends to the target, and standard patch tokenisation works against this by spreading capacity uniformly across the frame, starving a small object of representation precisely when it matters most. A foveated tokenisation that concentrates patches on the candidate region would impose a structural bias toward the target, and observing how that bias reshapes the description the model writes is a direct probe of how VLMs translate vision into language.

The second concerns how language is represented once written. We currently consume the text embeddings of the VLM and SAM~3 as they come, even though these representations are generic rather than tuned to correspondence, and a representation conditioned on the downstream task might carry the same phrase more usefully. This gain trades against the very property that motivates the paradigm, however, since a text embedding specialised far enough begins to resemble the opaque learned features we set out to avoid. How far language can be specialised before it stops being legible is therefore a question about the paradigm, not only about the pipeline.

The third is the most fundamental. Our results show that a single phrase reaches within a few IoU points of methods that pass continuous embeddings between their components, which raises the question of which medium compresses object identity more effectively. Language is a discrete channel of low bandwidth that a human can read, whereas a learned embedding carries far more and remains opaque. Measuring how much is lost in the compression to language, and whether that loss is fundamental or an artefact of current models, would tell us when the legible medium is worth its cost.

These directions build on a paradigm that already works rather than rescue one that failed. The broader claim the abstract opened therefore holds, that natural language can act as a reliable bridge between machine learning architectures, compressing what a learned latent space would otherwise carry into a phrase a human can read, question, and correct. The beauty of this paradigm lies in how it reconciles two goals that learned correspondence has long traded against each other, reaching the performance of specialised architectures while keeping every step of its reasoning legible. As foundation models continue to absorb a more general understanding of the visual and linguistic world, the most capable systems may increasingly be the ones we compose rather than the ones we train. Language may well remain the simplest and most human medium through which to compose them.
